\documentclass[pdf,notes]{beamer}
\mode<presentation>{\usetheme[secheader]{Boadilla}}

\usepackage{mystyle}
\includecomment{versiona}
\usepackage{xeCJK}
\usepackage{lpic}

\newcommand{\red}[1]{{\color[rgb]{0.6,0,0}#1}}

\makeatletter
\newenvironment<>{proofs}[1][\proofname]{\par\def\insertproofname{#1\@addpunct{.}}\usebeamertemplate{proof begin}#2}
{\usebeamertemplate{proof end}}
\makeatother

\makeatletter
\def\th@mystyle{%
	\normalfont % body font
	\setbeamercolor{block title example}{bg=orange,fg=white}
	\setbeamercolor{block body example}{bg=orange!20,fg=black}
	\def\inserttheoremblockenv{exampleblock}
}
\makeatother

\theoremstyle{mystyle}
\newtheorem{prop}{Proposition}

\setCJKmainfont[AutoFakeBold=true]{Hiragino Mincho Pro} %my Mac

\title{2つのゲーゲンバウアー多項式に関連する積分公式について\footnote{60 min talk}}

% A subtitle is optional and this may be deleted

\author[レオンチエフ・アレックス]{小林俊行\inst{1} \and \underline{レオンチエフ・アレックス}\inst{2}}

\institute[東大] % (optional, but mostly needed)
{
  \inst{1}%
  大学院数理科学研究科、カブリ数物連携宇宙研究機構\\
  東京大学
  \and
  \inst{2}%
  大学院数理科学研究科\\
  東京大学
  }

  \date[表現論とその周辺分野の広がり]{RIMS共同研究（公開型）\\「表現論とその周辺分野の広がり」\\京都大学数理解析研究所}
% - Either use conference name or its abbreviation.
% - Not really informative to the audience, more for people (including
%   yourself) who are reading the slides online

\subject{表現論}

\begin{document}
\begin{frame}\titlepage\end{frame}

\begin{frame}{Outline}
	\tableofcontents
\end{frame}
\section{section}
\begin{frame}
%%	\begin{figure}[h]
%%		\centering
%%		\includegraphics[scale=0.8]{wanted.eps}
%%		\label{fig:wanted}
%%	\end{figure}
\begin{lpic}[]{wanted(0.8)}
%%		\lbl[bl]{70,40;\scriptsize $\bullet$}
%%		\lbl[bl]{57,30;\scriptsize $ Y_{+,-\lambda}^{p,q}=\myabra{m-n\ge a+q}$}
%%		\lbl[bl]{25,70; $\myInd{\lambda}$}
%%
%%		\lbl[bl]{10,-11,45; \makebox[3cm]{\dotfill}}
%%		\lbl[bl]{9.4,-10,90; \makebox[3cm]{\hrulefill}}
%%		\lbl[bl]{9.0,-10; $\bullet$}
%%		\lbl[bl]{13.8,-11; \scriptsize$ 2-a-q =\lambda+\frac{p-q}{2}+1= b\left( \lambda,q,p \right)$}
%%		\lbl[bl]{71,19; \scriptsize$ =-\lambda-\frac{p-q}{2}-1=-b\left( \lambda,q,p \right)$}
		
%%		\lbl[bl]{25,70; \makebox[3cm]{\dotfill}} %dot 
%%		\lbl[bl]{25,65; \makebox[3cm]{\hrulefill}} %solid
%%		\lbl[bl]{25,65; \hdashrule[0.5ex]{3cm}{1pt}{3mm}} %dash
	\end{lpic}
\end{frame}
\end{document}


